% =====================================================================
% CONCLUSION
% =====================================================================
\section{نتیجه‌گیری}
\label{sec:conclusion}

ما به‌جای مدلِ درون یک برنامهٔ وبِ مبتنی بر \en{LLM}، خودِ برنامه را
ارزیابی کردیم، به شکلی که مهندس دیگری بتواند بازپخش و ممیزی‌اش کند، با
جدی‌گرفتنِ تعریف \en{Garak} از \en{generator} و بیانِ کلِ ارزیابی به‌شکل
\en{plugin}هایی که در زمان اجرا در فضای نام \en{Garak} ثبت می‌شوند. هستهٔ
\en{Garak} دست‌نخورده است، هر \num{28} اجرا از پیکربندی‌های نگه‌داشته‌شده
با \en{CLI} استاندارد بازپخش می‌شوند، و لایهٔ تحلیل ورودی‌های ارزیابی
\en{Garak} را تجمیع می‌کند بی‌آنکه هرگز پیامدی را دوباره قضاوت کند.

نتایج محتوایی نسبی‌اند، و عمداً چنین‌اند. \en{data poisoning}
پرمخاطره‌ترین دسته بود و کم‌وابسته‌ترین دسته به همکاریِ مدل. سخت‌کردن
\en{prompt} بهبودی واقعیِ \num{5.7} برابری داد که با این حال یک مرز نیست،
چون کنترلی که \num{$8.3\,\%$} مواقع در برابر مهاجمی شکست می‌خورد که بابت
هر تلاش چیزی نمی‌پردازد، زیر تکرار با قطعیت شکست می‌خورد. ثابت
نگه‌داشتنِ پیکرهٔ حمله در ده پیکربندی \en{guardrail}، شکست را به‌جای قدرتِ
یک دفاع به یک لایه نسبت داد و نشان داد دو مورد از ده مورد بهتر از نبودِ
\en{guardrail} عمل نمی‌کنند؛ قاعدهٔ نرمِ \en{system prompt} آموزنده است،
چون بیانِ قاعده‌ای دربارهٔ یک راز مستلزم قراردادنِ آن راز در \en{context}
است، پس آن قاعده هم‌زمان پیش‌شرط آسیب‌پذیری و نقشهٔ راه مهاجم است. موفقیت
در تکنیک‌هایی متمرکز شد که هرگز هدف خود را نام نمی‌برند، و تیزتر از همه
جایی که \en{prompt}‌هایی که \emph{قواعد} مدل را می‌خواستند رازِ
محافظت‌شده را \num{5.6} برابر بیشتر از \en{prompt}‌هایی استخراج کردند که
خودِ راز را می‌خواستند. پنج نظر با برچسب نادرست یک \en{backdoor} هدفمند در
طبقه‌بند ساختند که دقت تجمیعی را دست‌نخورده می‌گذارد و رانشش چهار گامِ
بودجه پیش از تغییر هر برچسبی قابل اندازه‌گیری است، و بازیابیِ
برچسب‌خوردهٔ اعتماد بازیابیِ غیرقابل‌اعتماد را کاملاً حذف کرد.

نتیجهٔ روش‌شناختی همان است که انتظار داریم دورتر از همه منتقل شود، و
\en{ablation} جایی است که به‌جای ادعا، فیصله می‌یابد. برداشتنِ کانال
\en{ground truth} دو \en{suite} بزرگ‌تر را به اندازه‌ای تضعیف نمی‌کند،
بلکه از میان برمی‌داردشان. جایگزینیِ یک \en{detector} عمومی استاندارد،
نقض‌های گزارش‌شده را با \en{recall} بدون تغییر \num{3.4} برابر می‌کند؛
مصرف‌کردنِ پرچمِ نشتِ خودِ برنامه به‌جای جدول‌بندیِ متقابلش، \num{0.762} را
جایی گزارش می‌کرد که \en{detector} مستقل \num{0.000} اندازه گرفت؛ و یک
جزئیاتِ چهارخطی در \en{harness} است که یک جاروب راستین را از همان
پیکربندی که $N$ بار اجرا شده و روندی روی آن گزارش شده جدا می‌کند. سه
مؤلفهٔ دیگر هیچ عددی را تغییر ندادند و با این حال در جدول هستند، چون
مطالعه‌ای که تنها کنترل‌هایی را گزارش کند که اتفاقاً شلیک کرده‌اند، دستگاهی
را توصیف می‌کند که به آن نیاز داشته نه دستگاهی که نتیجه را قابل اعتماد
می‌کند. همین انضباط بر تازه‌ترین افزودهٔ مجموعه هم حاکم است:
\en{LLM-as-a-judge} یک نظر دومِ اختیاری و کالیبره‌شده است که ترتیب
\en{schema} آن به‌جای فرض‌شدن اندازه‌گیری شده، و هیچ نرخ گزارش‌شده‌ای به آن
وابسته نیست.

با نگاه به همهٔ یافته‌ها، کنترل‌هایی که دوام آوردند روی خط لولهٔ برنامه عمل
می‌کنند و کنترل‌هایی که شکست خوردند روی سرشتِ مدل. نتیجهٔ فرعیِ آن برای
شدت، همان تک‌اصابتِ \en{SQL} میان‌مستأجری در \num{21} تلاش است: نرخ ناچیز
است و یافته قطعی، چون یک \en{routing flag} حدس‌ناپذیرِ مختصِ کاربر از راه
پرس‌وجویی که برنامه به دستور مدل اجرا کرد به فراخواننده رسید.

هر نرخ مطلقی اینجا کرانی پایینی برای یک استقرار پشت یک مدل یک‌میلیاردی
است. چهار مسیر تصویر را بی‌آنکه روش را تغییر دهند تیزتر می‌کنند: تکرار
همان \num{28} \en{task} روی نردبانی از مقیاس مدل، که دستگاهش پیشاپیش
پارامتری شده و شواهدِ وابستگیِ آسیب‌پذیریِ مسموم‌سازی به مقیاس آن را فوری
می‌کند~\cite[\en{Table~2}]{bowen2025scaling}؛ یک حملهٔ چندنوبتی که نخست
مقدمهٔ کاذبی را که \en{B6} به آن نیاز دارد برقرار کند؛ جایگزینیِ
امتیازدهیِ سم بر پایهٔ حضور واژه با یک بررسیِ ریشه‌داریِ خروجی؛ و انتقال
جدول‌های مشتق‌شده‌ای که تنها در این مقاله وجود دارند به پیمانهٔ تحلیل، چون
تا آن زمان همان‌ها تنها بخشی از این مطالعه‌اند که فقط با دست بازتولید
می‌شود.
